\documentclass{ca-capstone}
\capstonerefs{msu-ent,galbraith,knuth-taocp2,sagemath}
\begin{document}
% Generated by `make bench CAPSTONE=euclid-orders-primitive-roots` (taskgen.capstone report) from
% bench/results.json: reference time x slack on the reference machine.
% Do not edit: change bench/config.yaml and re-run the benchmark.
\newcommand{\ReqSlack}{10}
\newcommand{\ReqBitsMax}{64}
\newcommand{\ReqSeconds}{0{,}01}

\capstonetitle{Алгоритмы Евклида, порядки и первообразные корни}

\goal
Требуется реализовать набор связанных алгоритмов вычислительной теории чисел: НОД и обратные
элементы, КТО, порядки в мультипликативной группе поля и поиск первообразного корня.

\functionality
\begin{enumerate}
  \item Бинарный алгоритм вычисления НОД.
  \item Расширенный алгоритм Евклида и нахождение мультипликативно обратного элемента.
  \item Пакетное обращение списка элементов с явной пометкой нулей и необратимых элементов.
  \item Китайская теорема об остатках для попарно взаимно простых модулей и алгоритм Гарнера.
  \item Вычисление порядка элемента $\Zn{p}^*$ последовательным снижением степени и построением
    дерева произведений. Простое~$p$ и факторизация $p-1$ даны во входных данных.
  \item Поиск первообразного корня по простому модулю~$p$ с использованием заданной факторизации
    $p-1$. Любой первообразный корень считается правильным ответом.
  \item Промежуточные трассировки коэффициентов Безу, КТО/Гарнера, снижения порядка и проверки
    первообразного корня.
\end{enumerate}

\requirements
\begin{enumerate}
  \item Независимая проверка всех результатов: тождество Безу, обратимость, система сравнений,
    минимальность порядка и критерий первообразного корня.
  \item Корректные краевые случаи: знаки и нули для НОД, необратимые элементы, перестановки
    модулей в системе, кратные простые в факторизации $p-1$, собственные делители порядка.
  \item Примеры с числами длиной до \textbf{\ReqBitsMax~бит} должны выполняться не более чем за
    \textbf{\ReqSeconds~с} на запрос.
\end{enumerate}

\details
Разрешается использовать любой язык программирования, в частности Python, систему 
компьютерной алгебры SageMath и встроенные большие целые. Библиотечные \texttt{gcd},
\texttt{xgcd}, \texttt{inverse\_mod}, \texttt{crt}, \texttt{multiplicative\_order} и
\texttt{primitive\_root} в реализации запрещены и используются только как оракул. Факторизация
в работе не выполняется: разложение $p-1$ включается в набор входных данных.

Требуются поддержка только простых модулей в задачах вычисления порядка и поиска первообразных 
корней. Подъём на $p^k$, модули $2p^k$, рациональная реконструкция, общая КТО и перечисление всех
первообразных корней не входят в задание.

\schedule
Работа выполняется индивидуально. Срок выполнения -- \textbf{месяц}; ожидаемая трудоёмкость
-- примерно 38-42 часов. Реализация должна быть разработана и написана вами.

\remarks
\begin{enumerate}
  \item Пакетное обращение должно использовать один основной обратный элемент для обратимой
    части входа, а не скрыто вызывать обращение отдельно для каждого элемента.
  \item Векторы с нулями и необратимыми элементами являются обязательными.
  \item Оптимизация кода не требуется.
\end{enumerate}

\references
\end{document}
